\section{数据集与任务定义}

本章主要说明实验采用的数据资源、预处理流程、数据切分策略、待解决的问题以及评估指标。

\subsection{数据集介绍}

本实验选用 Thought2Text 数据集作为研究基础。该数据集采集自六名健康受试者，每位受试者依次观看隶属于 40 个 ImageNet 类别的自然图像，同时在观看期间以 128 通道记录其 EEG 信号。整个数据集包含约 11965 条 EEG trial，对应大约 1996 张不同的图像。原始 EEG 经降采样后变为 64 通道、250 个时序采样点，即维度为 $[64, 250]$。图像以 $224\times224$ 的分辨率输入模型。

需要指出的是，Thought2Text 本质上是一个 EEG-图像配对数据集，其原始标注只有类别标签，不包含自然语言描述。这一特性意味着该数据集更适合用来研究 EEG 与图像之间的类别级语义关联，而非直接训练 EEG-to-text 的端到端生成。为弥补这一不足，我们利用预训练的 BLIP-base 模型为每张图像自动生成伪描述文本，将其作为生成式模型的训练目标。

表~\ref{tab:dataset} 汇总了数据集的核心信息。

\begin{table}[htbp]
  \centering
  \caption{Thought2Text 数据集基本信息}
  \label{tab:dataset}
  \begin{tabular}{l c}
    \toprule
    项目 & 数值 \\
    \midrule
    图像总数 & 约 1996 张 \\
    EEG trial 总数 & 约 11965 条 \\
    类别数 & 40 个 \\
    被试数 & 6 名 \\
    EEG 形状 & $64\times250$ \\
    图像输入分辨率 & $224\times224$ \\
    原始 EEG 通道数 & 128 \\
    \bottomrule
  \end{tabular}
\end{table}

\subsection{数据预处理}

整个预处理流程涉及 EEG 信号、图像以及文本三个方面的处理。

\textbf{EEG 预处理}：原始 EEG 信号经过 128 到 64 通道的降采样，随后对各通道进行逐通道的 z-score 归一化。所有 EEG trial 被统一截取为 250 个采样点的长度。为了降低训练阶段的磁盘读写负担，本设计将训练集、验证集和测试集的所有 EEG 数据提前缓存为 numpy 数组格式，训练过程中直接从内存载入。

\textbf{图像处理}：所有图像均按 CLIP ViT-B/32 的标准预处理流程处理，包括缩放到 $224\times224$、中心裁剪以及 ImageNet 统计量归一化。为节省计算资源，我们预先使用 CLIP 视觉编码器提取全量图像的 pooled embedding 和 patch tokens，以 float16 精度保存为 .npy 文件，训练阶段读入后转换为 bfloat16 参与运算。

\textbf{文本处理}：针对数据集中全部的 40 个语义类别，本设计按模板 \texttt{"a photo of a \{class\_name\}"} 构造文本描述，经由冻结的 CLIP text encoder 编码为 512 维的文本原型向量。这些向量同样预先计算并存储，在训练和评估过程中保持不变。

\textbf{视觉退化生成}：六种视觉退化条件均在数据加载阶段在线施加，无需预先生成并保存退化后的图像。退化变换基于 torchvision 的标准算子实现，在 DataLoader 的 worker 进程内并行完成。

\subsection{数据划分}

本设计采用 image-level split（按图像划分）而非 trial-level split（按单条 trial 划分）的策略来分割数据集。这一选择的核心考量在于防止信息泄露：若按 trial 分割，同一张图像的不同 EEG trial 可能分布在训练集和测试集中，模型便有可能利用图像视觉特征的记忆来完成分类，而非真正学习 EEG 信号中的语义信息。采用 image-level split 能够确保训练集和测试集之间不存在任何共享图像。

各子集的具体规模如表~\ref{tab:split} 所示。其中验证集主要用于模型选择与早停判断，测试集仅在最终评估阶段使用。

\begin{table}[htbp]
  \centering
  \caption{数据集训练/验证/测试划分}
  \label{tab:split}
  \begin{tabular}{l c c c}
    \toprule
    划分 & 图像数 & EEG trial 数 & 类别数 \\
    \midrule
    训练集 & 1330 & 7970 & 40 \\
    验证集 & 333 & 1998 & 40 \\
    测试集 & 333 & 1997 & 40 \\
    \bottomrule
  \end{tabular}
\end{table}

\subsection{任务定义}

本实验聚焦于两个相互关联的核心任务。

\textbf{任务一：EEG 辅助语义预测}。以退化图像及其对应的 EEG trial 作为模型输入，输出在 40 个类别上的概率分布。该任务旨在检验真实配对 EEG 是否能够在视觉退化条件下，为模型判断图像类别提供额外信息。评估时涵盖五种输入设定：纯视觉（vision-only）、真实 EEG（real EEG）、shuffled EEG（打乱脑电）、random EEG（随机脑电）以及仅使用 EEG（eeg-only）。

\textbf{任务二：生成式图像描述}。以退化图像及其对应的 EEG trial 作为输入，模型需生成一段自然语言描述文本。该任务旨在验证脑电增强的视觉表示能否使生成式视觉语言模型产生比纯视觉输入更准确的描述。评估时包含四种设定：纯视觉（vision-only）、真实 EEG（real EEG）、shuffled EEG（打乱脑电）以及random EEG（随机脑电）。

\subsection{评价指标}

\textbf{语义预测指标}：

\begin{itemize}
  \item \textbf{Top-1 Accuracy}：模型预测概率最高的类别与真实类别一致的样本比例；
  \item \textbf{Top-5 Accuracy}：真实类别出现在模型预测概率前五位中的样本比例；
  \item \textbf{Class Hit}：利用 CLIP 文本原型相似度进行零样本类别判定，命中真实类别的比率；
  \item \textbf{Real-Vision Gap}：真实 EEG 模式与纯视觉模式的指标差值，正值表明 EEG 带来了正向增益；
  \item \textbf{Real-Shuffled/Random Gap}：真实 EEG 与打乱脑电/随机脑电的指标差值，是判断 EEG 效果是否源自真实配对关系的关键对照指标；
  \item \textbf{Win Rate}：在多种退化条件或随机种子中，真实 EEG 模式优于对照组的实验比例。
\end{itemize}

\textbf{生成式指标}：

\begin{itemize}
  \item \textbf{有效输出率（Valid Caption Rate）}：生成文本中合法有效描述所占的比例。有效描述定义为非空、非纯乱码、非代码片段且非 URL 的文本；
  \item \textbf{无效输出率（Invalid Output Rate）}：生成文本中无效输出的比例；
  \item \textbf{Caption Class-Hit}：对生成描述进行 CLIP 零样本类别判定，命中真实类别的比率；
  \item \textbf{Distinct Caption Count}：生成的不同描述文本的总数，反映输出结果的多样性；
  \item \textbf{平均描述长度（Avg Caption Length）}：生成描述的平均 token 数目。
\end{itemize}
